\chapter{\objectivesname}
Este Trabajo Fin De Grado tiene como objetivo ofrecer una visión actualizada mediante el análisis de manera sistemática de la literatura científica de las comunicaciones y ciberseguridad en la era post-cuántica.

El estudio busca considerar mas allá de aspectos teóricos, la implementación práctica en escenarios reales, respondiendo a la necesidad de contextualizar los avances tecnológicos cuánticos dentro de entornos aplicables, permitiendo identificar limitaciones tecnológicas, estrategias de migración, propuestas de líneas futuras que puedan orientar el despliege seguro y efectivo de estas tecnologías dentro de la industria y sistemas críticos.

Con el fin de estructurar el alcance de este estudio, se establecen los siguientes objetivos:

\begin{itemize} % Lista de objetivos específicos
    \item Analizar tecnologías habilitadoras de comunicaciones cuánticas, incluyendo distribución cuántica de claves (QKD), repetidores cuánticos, canales cuánticos y protocolos de transmisión, evaluando su desarrollo actual y limitaciones.
    \item Evaluar las arquitecturas e infraestructuras de redes cuánticas, incluyendo redes híbridas de fibra óptica, satelitales, considerando su escalabilidad, viabilidad práctica e integración con las infraestructuras existentes.
    \item Examinar los mecanismos de seguridad y criptografía cuántica y post-cuántica, considerando la adaptación de protocolos tradicionales, estrategias de migración y protección frente a nuevas vulnerabilidades introducidas por la computación cuántica.
    \item Identificar amenazas emergentes potenciadas por la computación cuántica y los mecanismos de defensa y resiliencia frente a ellas.
\end{itemize}

En su conjunto, este estudio no solo busca documentar el estado del arte, sino que aspira a ser una herramienta de recopilación de todo el progreso en tecnologías cuánticas y post-cuánticas dentro de la seguridad y sistemas de comunicación, funcionando además como un punto de referencia que facilita el acceso a información clave y proporciona una base sólida para futuros proyectos de investigación académica. 


